\documentclass[12pt]{article}
\usepackage{fancyhdr}
\pagestyle{fancy}
\fancyhf{}
\fancyfoot[C]{\thepage}
\fancyfoot[L]{\scriptsize CC BY-NC-SA 4.0 \textbullet\ Leo Liang}
\renewcommand{\headrulewidth}{0pt}
\usepackage[margin=1in]{geometry}
\usepackage{amsmath, amssymb}
\usepackage{array, booktabs}
\usepackage{pgffor}
\usepackage{graphicx}
\parindent=0pt
\parskip=1em
\newcolumntype{C}[1]{>{\centering\arraybackslash}p{#1}}
% writing lines
\newcommand{\wl}{\par\vspace{5mm}\noindent\rule{\linewidth}{0.4pt}\par\vspace{1mm}}
\newcommand{\wls}[1]{\foreach \i in {1,...,#1}{\wl}}
% ─────────────────────────────────────────────────────────────────────────────
\begin{document}
\pagestyle{fancy}
\begin{center}
  {\Large\bfseries Homework 3}\\[4pt]
  Introduction to Functions\\[6pt]
  Name:\;\underline{\hspace{6.5cm}}\quad
\end{center}
\hrule\bigskip
% =============================================================================
\section*{Instructions}
Use a \textbf{pen}. Write questions and surprises in the margin in a different color.
Some statements in this homework are intentionally imprecise, and some contradict what you already
know --- find them and flag them in the margin. If a topic confuses you, you can use the internet or AI to
help. AI can be a good teacher if you ask it specific questions. \textbf{DO NOT} use AI to solve the problems
without trying yourself first.

Remember to write comments in the margin and answer the entire question. You are a mathematician, your answers should be detailed and convincing to anyone that reads it. \textbf{Do not leave ambiguity} (remember what ambiguous means? if not, search it up).

% =============================================================================
\section*{Part 1 --- Foundation}

\textbf{1.}\quad Let $f(x) = 3x + 4$. Fill in the table, then answer the questions.

\medskip
\renewcommand{\arraystretch}{2.4}
\begin{tabular}{|C{4cm}|C{4cm}|}
\hline
$x$ & $f(x)$ \\
\hline
$0$ & \\
\hline
$1$ & \\
\hline
$-1$ & \\
\hline
$2$ & \\
\hline
$\tfrac{1}{3}$ & \\
\hline
\end{tabular}
\renewcommand{\arraystretch}{1}

\medskip
\textbf{(a)}\quad What is $f(f(1))$? \hfill Answer:\;\underline{\hspace{3cm}}

\medskip
\textbf{(b)}\quad Find $x$ such that $f(x) = 10$. \hfill Answer:\;\underline{\hspace{3cm}}

\medskip
\textbf{(c)}\quad Find $x$ such that $f(x) = 0$. \hfill Answer:\;\underline{\hspace{3cm}}

\bigskip
\textbf{2.}\quad For each rule, decide whether it defines a valid function from $A$ to $B$. If it
does, state whether it is injective and/or surjective. If it does not, explain why not.

\medskip
\textbf{(a)}\quad $A = \{1,2,3\}$, $B = \{1,2,3,4,5\}$:\quad $1 \mapsto 2$, $2 \mapsto 4$, $3 \mapsto 5$.
\wls{2}

\textbf{(b)}\quad $A = \{1,2,3\}$, $B = \{1,2,3,4,5\}$:\quad $1 \mapsto 3$, $2 \mapsto 3$, $3 \mapsto 3$.
\wls{2}

\textbf{(c)}\quad $A = \{1,2,3\}$, $B = \{1,2,3,4,5\}$:\quad $1 \mapsto 2$, $2 \mapsto 4$, $2 \mapsto 5$.
\wls{2}

\textbf{(d)}\quad $A = \{1,2,3\}$, $B = \{1,2,3\}$:\quad each number maps to one more than itself,
$x \mapsto x+1$.
\wls{2}

\textbf{(e)}\quad $A = \{1,2,3\}$, $B = \{a,b,c\}$:\quad $1 \mapsto a$, $2 \mapsto b$.
\wls{2}

\bigskip
\textbf{3.}\quad Domain, codomain, and range.

\medskip
\textbf{(a)}\quad $g\colon \mathbb{R} \to \mathbb{R}$ defined by $g(x) = x^2 + 2$.
State the domain, codomain, and range.
\wls{2}

\textbf{(b)}\quad $h\colon \{1,2,3\} \to \{3,6,11\}$ defined by $h(x) = x^2 + 2$.
State the domain, codomain, and range. Is $h$ surjective? You used the same formula as (a) ---
explain why the answers differ.
\wls{3}

\textbf{(c)}\quad \textit{No formula required.}\quad A Minecraft furnace smelts each input into one output:
\[
  \text{smelt}\colon \{\text{iron ore},\,\text{gold ore},\,\text{sand},\,\text{raw beef}\}
  \longrightarrow \{\text{iron ingot},\,\text{gold ingot},\,\text{glass},\,\text{steak}\},
\]
\[
  \text{iron ore}\mapsto\text{iron ingot},\;\;
  \text{gold ore}\mapsto\text{gold ingot},\;\;
  \text{sand}\mapsto\text{glass},\;\;
  \text{raw beef}\mapsto\text{steak}.
\]
State the domain, codomain, and range of smelt. Can you write smelt as a formula $\text{smelt}(x) = \ldots$?
What does your answer say about the claim ``every function is a formula''?
\wls{3}

\bigskip
\textbf{4.}\quad Fill in the blanks, then answer the True/False items (justify each in one sentence).

\medskip
\textbf{(a)}\quad The range of a function is always a \underline{\hspace{3cm}} of its codomain.

\textbf{(b)}\quad In symbols, for $f\colon A \to B$, the range $R(f)$ satisfies $R(f) \subseteq \underline{\hspace{2.5cm}}$.

\textbf{(c)}\quad A function is surjective exactly when its range equals its \underline{\hspace{3cm}}.

\medskip
T\,/\,F\quad For every function, the range is a subset of the codomain.
\wl

T\,/\,F\quad For every function, the range equals the codomain.
\wl

T\,/\,F\quad Every function is surjective onto its own range (shrink the codomain to the range and it becomes surjective).
\wl

\textbf{(d)}\quad Give a function $f\colon \{1,2\} \to \{a,b,c\}$ whose range is a \textbf{proper} subset of its codomain. Write it as ordered pairs, and state the range.
\wls{2}

\bigskip
\textbf{5.}\quad For finite sets, the sizes $|A|$ and $|B|$ alone determine what kinds of functions can exist. Answer True/False and justify each in one sentence.

\medskip
T\,/\,F\quad If $|A| > |B|$, then no function $f\colon A \to B$ can be injective.
\wl

T\,/\,F\quad It is possible for $f\colon A \to B$ to be injective when $|A| \ge |B|$.
\wl

T\,/\,F\quad If $|A| < |B|$, then no function $f\colon A \to B$ can be surjective.
\wl

T\,/\,F\quad If $|A| = |B|$ (finite) and $f$ is injective, then $f$ must also be surjective.
\wl

T\,/\,F\quad If $f\colon A \to B$ is bijective, then $|A| = |B|$.
\wl

\medskip
\textbf{(a)}\quad Complete the table. For each pair of sizes, circle whether each type of function from $A$ to $B$ is \textbf{possible}.

\medskip
\renewcommand{\arraystretch}{2.2}
\begin{tabular}{|C{3.2cm}|C{3.2cm}|C{3.2cm}|C{3.2cm}|}
\hline
Sizes & Injective possible? & Surjective possible? & Bijective possible? \\
\hline
$|A|=2,\ |B|=4$ & Yes \quad No & Yes \quad No & Yes \quad No \\
\hline
$|A|=4,\ |B|=2$ & Yes \quad No & Yes \quad No & Yes \quad No \\
\hline
$|A|=3,\ |B|=3$ & Yes \quad No & Yes \quad No & Yes \quad No \\
\hline
\end{tabular}
\renewcommand{\arraystretch}{1}

\bigskip
\textbf{6.}\quad \textbf{One Function, Four Outfits.}\quad
A function can be dressed four ways --- an \textbf{arrow diagram}, a set of \textbf{ordered pairs}, a \textbf{table}, and a \textbf{graph}. Here is one function, given as ordered pairs:
\[
  f\colon \{-1,0,1,2\} \to \mathbb{Z}, \qquad f = \{(-1,\,1),\,(0,-1),\,(1,-3),\,(2,-5)\}.
\]

\textbf{(a)}\quad Complete the table.

\medskip
\renewcommand{\arraystretch}{2.4}
\begin{tabular}{|C{2.4cm}|C{2.2cm}|C{2.2cm}|C{2.2cm}|C{2.2cm}|}
\hline
$x$ & $-1$ & $0$ & $1$ & $2$ \\
\hline
$f(x)$ & & & & \\
\hline
\end{tabular}
\renewcommand{\arraystretch}{1}

\newpage

\medskip
\textbf{(b)}\quad Draw the arrow diagram, and plot the points on a grid.

\medskip
\begin{tabular}{|p{0.46\linewidth}|p{0.46\linewidth}|}
\hline
\textbf{Arrow diagram} & \textbf{Graph} \\
& \\[2.8cm]
\hline
\end{tabular}

\medskip
\textbf{(c)}\quad Find the \textbf{rule}. $f(x) = \underline{\hspace{4cm}}$

\medskip
\textbf{(d)}\quad Is $f$ injective? Is $f$ surjective onto $\mathbb{Z}$? For each, name the representation that made it easiest to see, and say why.
\wls{3}

\bigskip
\textbf{7.}\quad \textbf{Is It a Function? Reading the Graph.}\quad
The \textbf{vertical line test}: a graph in the $xy$-plane defines $y$ as a function of $x$ exactly when no vertical line meets it more than once. For each, decide whether it defines $y$ as a function of $x$ --- that is, whether there is a function $f$ with $y = f(x)$. Circle Yes / No and justify in one sentence.

\medskip
\textbf{(a)}\quad The points $\{(1,2),\,(2,4),\,(1,5),\,(3,4)\}$. \hfill Yes\quad No
\wl

\textbf{(b)}\quad The horizontal line $y = 3$. \hfill Yes\quad No
\wl

\textbf{(c)}\quad The vertical line $x = 3$. \hfill Yes\quad No
\wl

\textbf{(d)}\quad The parabola $y = x^2 - 1$. \hfill Yes\quad No
\wl

\textbf{(e)}\quad The sideways parabola $x = y^2$. \hfill Yes\quad No
\wl

\textbf{(f)}\quad The circle $x^2 + y^2 = 16$. \hfill Yes\quad No
\wl

\textbf{(g)}\quad In your own words, why does the vertical line test work? (What does a vertical line hitting the graph twice say about a single input?)
\wls{2}

\bigskip
\textbf{Reading --- The Two Arrows: $\to$ and $\mapsto$}

\medskip
Function statements use two different arrows, and they are \textbf{not} interchangeable (do you know what interchangeable means?).

\medskip
\textbf{The long arrow $\to$ joins two \emph{sets}.}\quad Writing $f\colon A \to B$ says ``$f$ is a function from the set $A$ to the set $B$.'' It names the domain and the codomain. Read it ``$f$, from $A$ to $B$.''

\medskip
\textbf{The bar arrow $\mapsto$ joins two \emph{elements}.}\quad Writing $1 \mapsto a$ says ``the input $1$ is sent to the output $a$'' --- the very same statement as $f(1) = a$. Read it ``$1$ maps to $a$.'' The little bar at the tail reminds you it starts at one specific element.

\medskip
So one fact can be written three equivalent ways:
\[
  f(1) = a \qquad\Longleftrightarrow\qquad 1 \mapsto a \qquad\Longleftrightarrow\qquad (1,a) \in f.
\]

Putting it together, a line like
\[
  p\colon \{1,2,3\} \to \{a,b,c,d\}, \qquad 1 \mapsto a,\;\; 2 \mapsto b,\;\; 3 \mapsto c
\]
first declares the domain and codomain (the $\to$ part), then lists where each input goes (the $\mapsto$ part). It is the same function as the ordered-pair set $\{(1,a),(2,b),(3,c)\}$.

\bigskip
\textbf{8.}\quad \textbf{Injective, Surjective, Bijective --- Spot Them, Then Build Them.}

\medskip
\textbf{Part 1 --- Label} each function as \textbf{injective}, \textbf{surjective}, \textbf{bijective}, or \textbf{none}.

\medskip
\textbf{(a)}\quad $p\colon \{1,2,3\}\to\{a,b,c,d\}$:\quad $1\mapsto a,\;2\mapsto b,\;3\mapsto c$.
\wl

\textbf{(b)}\quad $q\colon \{1,2,3\}\to\{a,b\}$:\quad $1\mapsto a,\;2\mapsto b,\;3\mapsto a$.
\wl

\textbf{(c)}\quad $r\colon \{1,2,3\}\to\{a,b,c\}$:\quad $1\mapsto c,\;2\mapsto a,\;3\mapsto b$.
\wl

\textbf{(d)}\quad $s\colon \{1,2,3\}\to\{a,b,c\}$:\quad $1\mapsto a,\;2\mapsto a,\;3\mapsto b$.
\wl

\medskip
\textbf{Part 2 --- Build your own} (write each as a set of ordered pairs).

\medskip
\textbf{(e)}\quad An \textbf{injective but not surjective} function $\{1,2\}\to\{a,b,c\}$.
\wl

\textbf{(f)}\quad A \textbf{surjective but not injective} function $\{1,2,3\}\to\{a,b\}$.
\wl

\textbf{(g)}\quad A \textbf{bijective} function $\{1,2,3\}\to\{x,y,z\}$.
\wl

\newpage
\textbf{(h)}\quad $\bigstar$\quad Try to build a \textbf{bijective} function $\{1,2,3\}\to\{a,b\}$. Explain exactly what goes wrong.
\wls{2}

% =============================================================================
\section*{Part 2 --- Application}

\textbf{1.}\quad \textbf{Three Bedrooms.}\quad
Back to the Minecraft build. You are now carpeting \textbf{three} bedrooms. Every room is $10$ blocks long.
Bedroom~1 is $6$ blocks wide and Bedroom~2 is $8$ blocks wide --- both finished. Bedroom~3 has an
undecided width $w$. Each floor block needs one carpet block, and each carpet block takes $2$ wool.

\medskip
\textbf{(a)}\quad Write a function $f(w)$ giving the \textbf{total} wool needed for all three rooms.
\wls{2}

\textbf{(b)}\quad Evaluate $f(5)$. What does this number mean in context?
\wls{2}

\textbf{(c)}\quad You have collected exactly $480$ wool and want to use all of it. Find $w$ such that
$f(w) = 480$. What does this answer mean in context?
\wls{2}

\newpage
\textbf{(d)}\quad State the natural domain of $f$ in the Minecraft context. Is $f$ surjective onto the positive
integers $\mathbb{Z}^+$? Explain, and give one positive integer that is never an output.
\wls{3}

\bigskip
\textbf{2.}\quad \textbf{Speaking Functions.}

\medskip
\textbf{(a)}\quad Read $f(x) = 2x^2 - 5$ aloud as one complete English sentence. Write what you would say.
\wls{2}

\textbf{(b)}\quad A classmate says: ``$f(x)$ means $f$ times $x$.'' Are they right? Correct them in one sentence.
\wl

\textbf{(c)}\quad Read $f(g(x))$ aloud as a complete English sentence. What do you think this expression means?
\wls{2}

\newpage
\bigskip
\textbf{3.}\quad \textbf{Functions Live Inside $A \times B$.}\quad
Last week you learned the Cartesian product $A \times B$. A function $f\colon A \to B$ is a special subset of $A \times B$ in which each
element of $A$ appears as a first coordinate exactly once.

Let $A = \{1,2\}$ and $B = \{p,q,r\}$.

\medskip
\textbf{(a)}\quad List every element of $A \times B$.
\wl

\textbf{(b)}\quad How many ordered pairs are in $A \times B$? (Recall $|A \times B| = |A|\cdot|B|$.)
\wl

\textbf{(c)}\quad For each subset, circle whether it is a function from $A$ to $B$. For each that is not, say why.

\medskip
\hspace{1.5em}(i)\quad $\{(1,p),(2,r)\}$ \hfill function?\quad Yes\quad No
\wl

\hspace{1.5em}(ii)\quad $\{(1,p),(1,q)\}$ \hfill function?\quad Yes\quad No
\wl

\hspace{1.5em}(iii)\quad $\{(1,p)\}$ \hfill function?\quad Yes\quad No
\wl

\textbf{(d)}\quad $\bigstar$\quad How many functions are there from $A=\{1,2\}$ to $B=\{p,q,r\}$ in total?
Explain your count --- do not just guess.
\wls{2}

\newpage
% =============================================================================
\section*{Looking Ahead --- Why Functions, and Where We Go Next}
\textit{Read this. Some of it is meant to feel unfinished --- that is the point. It also contains at least
one statement that is not quite right; find it and flag it in the margin. Where you get stuck, write the
question down and we will resolve it next class.}

\medskip
Here is the story so far. We started with \textbf{sets} (collections of objects), built \textbf{ordered pairs}
and the \textbf{Cartesian product} $A \times B$, and discovered that a \textbf{function} is just a special
subset of $A \times B$. Then we learned to measure functions with \textbf{domain}, \textbf{codomain}, and
\textbf{range}, and to classify them as \textbf{injective}, \textbf{surjective}, or \textbf{bijective}. Next we
ask three new questions: how do functions \textit{combine}, how do we \textit{undo} them, and how can they be
used to \textit{measure infinity}?

\subsection*{1. Combining Functions --- Composition}
Think of two machines in a line. A furnace turns \textbf{ore} into an \textbf{ingot}; a crafting table turns an
\textbf{ingot} into a \textbf{tool}. Feed ore in one end and a tool comes out the other --- you have chained two
functions. If $f$ is ``smelt'' and $g$ is ``craft,'' the combined machine is the \textbf{composition}
$g \circ f$, read ``$g$ of $f$,'' where $(g \circ f)(x) = g(f(x))$.

\medskip
\textbf{Try it.}\quad Let $f(x) = x + 1$ and $g(x) = 2x$. Compute $(g \circ f)(3)$, then $(f \circ g)(3)$.
Did you get the same answer both ways?
\wls{2}

\subsection*{2. Undoing Functions --- Inverses}
In the Three Bedrooms problem you ran the wool function \textit{backwards}: you knew the output ($480$ wool) and recovered the
input ($w$). That ``backwards machine'' is called the \textbf{inverse}, written $f^{-1}$. An inverse takes each
output back to the input it came from. Because every function can be undone, every function has an inverse.

\medskip
\textbf{Try it.}\quad For $f(w) = 20w + 280$, describe in words the steps of the backwards machine that turns a
wool count back into a width.
\wls{2}

\medskip
\textbf{Answer this.}\quad What if two different inputs gave the \textit{same} output? Could the
backwards machine still know which input to return to? Explain.
\wls{3}

% =============================================================================
\section*{Reading --- Proofs and Why a Picture Is Not Enough}

\textit{Read this before Part 3. The Challenge problems ask you to \textbf{prove} things. So first, what is a proof, and why do mathematicians insist on them instead of just drawing a picture?}

\subsection*{A claim about a function is a claim about \emph{every} input}
When we say ``$f$ is injective'' we are promising something about \textbf{every} pair of inputs: no matter which two different inputs you pick, their outputs differ. When we say ``$f$ is surjective'' we promise that \textbf{every} element of the codomain gets hit. These are ``for all'' statements --- they sweep over the whole domain or codomain, which may be infinite.

An arrow diagram, a table, or a sketched graph shows you only the cases you actually drew. For a function on a tiny finite set like $\{1,2,3\}$ you can draw \textbf{all} the arrows, and then the picture really does check every case --- that is an honest proof by exhaustion, because you looked at everything. But the moment the domain is large or infinite --- $\mathbb{N}$, $\mathbb{Z}$, $\mathbb{R}$ --- you can never draw all the arrows or plot the whole graph. The picture becomes a \textbf{sample}, and a sample cannot settle a claim about infinitely many inputs.

\subsection*{Pictures can fool you}

A picture is where \textbf{ideas} come from --- it shows you what to believe and where to hunt for trouble. A \textbf{proof} is what makes you \textbf{certain}: a chain of logical steps, built from \textbf{definitions}, that forces the conclusion for every case at once.

\subsection*{The engine: reason about an \emph{arbitrary} element}
Here is the trick that lets a few lines of writing cover infinitely many cases. Instead of testing specific inputs, we say ``let $a$ and $b$ be \textbf{arbitrary}'' --- we assume nothing about which elements they are. Whatever we then prove about $a$ and $b$ must hold for \textbf{all} of them, because we never used anything special about them. One argument, every case. This is the whole reason a proof can beat infinitely many pictures.

\textbf{A note on the $\square$.}\quad The little box at the end of a proof is a full stop for arguments: it announces ``the proof is finished. Everything we promised has now been shown.'' It is shorthand for \textbf{QED}, the Latin \textit{quod erat demonstrandum}, meaning ``which was to be demonstrated.'' Mathematicians once spelled out the letters Q.E.D.; today most use the hollow box $\square$ (or a filled-in $\blacksquare$), a habit popularized by the mathematician Paul Halmos, who borrowed it from magazine typography. The symbol carries no mathematical content --- it is purely a sign to the reader, marking exactly where the reasoning ends so that nothing after it gets mistaken for part of the argument.

\subsection*{Proving injectivity}
\textbf{Definition.}\quad $f$ is injective if $f(a) = f(b)$ forces $a = b$ (recall this is the contrapositive of the definition of injectivity).

\textbf{Blueprint.}\quad We approach this by first supposing $f(a) = f(b)$ for arbitrary $a,b$ in the domain. Then we manipulate using the rule of $f$. Then if we can conclude $a = b$, we have proved $f$ is injective.

\medskip
\textbf{Worked example.}\quad Let $f\colon \mathbb{R} \to \mathbb{R}$, $f(x) = 3x + 4$.

Claim: $f$ is injective.

\textbf{Proof.}\quad Suppose $f(a) = f(b)$. Then $3a + 4 = 3b + 4$, so $3a = 3b$, therefore $a = b$. Since $a,b$ were arbitrary, $f$ is injective. $\square$

\medskip
Notice what no picture could do: this single argument covered every one of the infinitely many pairs of real numbers.

\medskip
\textbf{Disproving injectivity takes only ONE counterexample.}\quad To show $g(x) = x^2$ is \textbf{not} injective we need no general argument --- just one clash: $g(2) = 4 = g(-2)$, yet $2 \ne -2$. A single counterexample destroys a ``for all'' claim.

\subsection*{Proving surjectivity}
\textbf{Definition.}\quad $f\colon A \to B$ is surjective if for every $b \in B$ there exists some $a \in A$ with $f(a) = b$.

\textbf{Template.}\quad Take an \textit{arbitrary} target $y \in B$. \textit{Produce} an input $x \in A$ (usually by solving $f(x) = y$) and check that $f(x) = y$.

\medskip
\textbf{Worked example.}\quad Let $f\colon \mathbb{R} \to \mathbb{R}$, $f(x) = 3x + 4$.

\textit{Claim: $f$ is surjective.}

\textbf{Proof.}\quad Let $y \in \mathbb{R}$ be arbitrary. Set $x = \dfrac{y - 4}{3}$, which is a real number, so it lies in the domain. Then
\[
  f(x) = 3 \cdot \frac{y - 4}{3} + 4 = (y - 4) + 4 = y.
\]
So every target $y$ has a corresponding element in the domain. Therefore $f$ is surjective. $\square$

You may want to re-read the example above several times and convince yourself that it is true (ask yourself questions like ``why did we need to show $f(x) = y$?'')

\medskip
\textbf{Surjectivity depends on the codomain.}\quad Suppose we defined $g\colon \mathbb{R} \to \mathbb{R}$, $g(x) = x^2$. Then $g$ is \textbf{not} surjective because the target $y = -1$ has no real corresponding input $x$ with $x^2 = -1$. So $-1$ is never hit (one unreachable target is enough to disprove ``onto''). If we shrunk the codomain to $[0,\infty)$ then our function $g$ becomes surjective. A surjectivity proof is always a statement about the \textbf{declared} codomain.

\subsection*{Proving bijectivity}
\textbf{Definition.}\quad $f$ is bijective if it is both injective and surjective.

So a bijectivity proof is simply \textbf{two} proofs stacked: prove injective, then prove surjective. Our $f(x) = 3x + 4$ is both, so it is a bijection --- and that is exactly the condition under which $f$ has an inverse (the ``undo'' machine from the preview).

\newpage
% =============================================================================
\section*{Part 3 --- Challenge}

\textbf{1.}\quad \textbf{Proving the Wool Function.}\quad
In Part~2 you built the three-bedroom wool function $f(w) = 20w + 280$. Take both the domain and codomain to be the positive integers $\mathbb{Z}^+$ (whole blocks of width, whole counts of wool).

\medskip
\textbf{(a)}\quad Is $f$ injective? Prove it.
\wls{4}

\textbf{(b)}\quad Is $f$ surjective? Prove or disprove.
\wls{4}

\textbf{(c)}\quad Is $f$ bijective? Prove or disprove.
\wls{2}

\newpage
\textbf{(d)}\quad In one sentence, say what injectivity means \textbf{in the Minecraft context}: what does it guarantee about the relationship between a width and a wool count?
\wls{2}

\bigskip
\textbf{2.}\quad \textbf{Odd One Out.}\quad
Consider $f\colon \mathbb{Z} \to \mathbb{Z}$ defined by $f(n) = 2n + 1$.

\medskip
\textbf{(a)}\quad Is $f$ injective? Prove it.
\wls{3}

\textbf{(b)}\quad Is $f$ surjective? Prove or disprove. Remember that finding a counterexample counts as a disproof.
\wls{3}

\textbf{(c)}\quad Describe the range of $f$ in words. If you kept the same rule but shrank the codomain to make $f$ surjective, what set would the new codomain be?
\wls{2}

\newpage
\bigskip
\textbf{3.}\quad \textbf{The Finite Trap.}\quad
Try to build a function $f\colon \{1,2,3,4\} \to \{1,2,3,4\}$ that is injective but \textbf{not} surjective. Explain what happened and why.
\wls{4}

\bigskip
\textbf{4.}\quad \textbf{The Collatz Function.}\quad
Define $C\colon \mathbb{N} \to \mathbb{N}$ by
\[
  C(n) =
  \begin{cases}
    n/2 & \text{if } n \text{ is even} \\
    3n + 1 & \text{if } n \text{ is odd.}
  \end{cases}
\]
Starting from any positive integer, apply $C$ over and over.

\medskip
\textbf{(a)}\quad Compute the sequence starting at $n = 7$. Stop when you reach $1$. You can use a calculator.
\wls{3}

\textbf{(b)}\quad Compute the first $10$ terms starting at $n = 27$. (Just list them; you do not need to reach $1$.)
\wls{2}

\newpage
\textbf{(c)}\quad Is $C$ injective? Either prove it or find a counterexample. (Hint: can an even number and an
odd number ever land on the same output?)
\wls{3}

\textbf{(d)}\quad Write a \textbf{conjecture} about what always seems to happen when you apply $C$ repeatedly,
starting from any positive integer.
\wls{2}

\bigskip
\textbf{5.}\quad $\bigstar$\quad \textbf{As Many Evens as Whole Numbers?}\quad
Let $E = \{2,4,6,8,\ldots\}$ be the set of positive even numbers. Define the function $f\colon \mathbb{N} \to E$ by
$f(n) = 2n$.

\medskip
\textbf{(a)}\quad Is $f$ injective? Prove it.
\wls{4}

\newpage
\textbf{(b)}\quad Is $f$ surjective onto $E$? Prove it.
\wls{4}

\textbf{(c)}\quad We will say that a bijection between two sets means they have the same size. But $E$ is only \textit{part} of
$\mathbb{N}$. What strange thing does this say about the size of $\mathbb{N}$ compared to $E$?
\wls{5}

\newpage
\medskip
\textbf{(d)}\quad $\bigstar$\quad For \textbf{finite} sets, a proper subset always has strictly fewer
elements --- so if $B \subset A$ you would expect no way to match $A$ one-to-one into $B$. In parts (a)--(b) you matched $\mathbb{N}$ one-to-one with $E$, the even numbers, even though $E$ is a proper subset of $\mathbb{N}$ (it leaves out every odd number). What does this mapping show about the rule ``a proper subset is strictly smaller'' for infinite sets? What does it suggest about judging the size
of an infinite set?
\wls{5}

% =============================================================================
\section*{Part 4 --- Notebook Artifact}

\bigskip
\textbf{1.}\quad If you used the internet or AI for help, describe how you used it:
\wls{5}

% =============================================================================
% =============================================================================
% =============================================================================
\end{document}
